\chapter{The geometric slope inequality}

Let $S$ be a proper adelic curve over $K$, with place space $\Omega$,
sigma-algebra $\mathcal A$, measure $\nu$, and absolute values $\phi$.
We assume that the Archimedean completions carry their usual real or complex
absolute values, and that either $\mathcal A$ is discrete or $S$ has a common
countable subfield dense in every completion. In the application below, $K$ is
countable of characteristic zero, so the latter condition holds with the whole
field $K$. Write $K_\omega$ for the completion at $\omega$ and
$m_\infty=\nu(\Omega_\infty)<\infty$.

All bundles are finite-dimensional, with their actual completed norms.
Subspaces carry induced norms, quotients carry quotient norms, and duals carry
operator norms. Tensor products of Hermitian bundles carry Hilbert norms at
Archimedean places and epsilon norms at all other places.

For a positive-rank bundle $\overline A$, put
$\widehat\mu(\overline A)=\widehat{\deg}(\overline A)/\dim_K A$.
Mean-slope semistability means that every nonzero subspace with its induced
norms has slope at most $\widehat\mu(\overline A)$.
The extreme slopes $\widehat\mu_{\max}$ and $\widehat\mu_{\min}$ are,
respectively, the supremum of slopes of nonzero subspaces and the infimum of
slopes of nonzero quotients. These definitions use the extended real line when
necessary. For a semistable bundle the maximal slope is its mean slope; for a
Hermitian semistable bundle the minimal slope is also its mean slope.
We use the degree, tensor, and duality formulas recalled in
\cite[Section 2.4]{bost2026}.

\section{Contractions and sublinear errors}

\begin{lemma}[Semistable slopes under a nonzero contraction]
\label{lem:semistable-contraction-slope}
\uses{def:adelic-slope,lem:slope-duality}
\source{bost-adelic-2026:page-0006,bost-adelic-2026:page-0007,bost-adelic-2026:page-0027}
\lean{ArakelovGeometry.AdelicVectorBundle.slope_le_of_nonzero_completed_contraction}
\leanok
Let $\overline A$ and $\overline B$ be mean-slope semistable adelic bundles,
and suppose that $\overline A$ is Hermitian. If a nonzero $K$-linear map
$f:A\to B$ satisfies
\[
 \|f_\omega(x)\|_{B,\omega}\leq\|x\|_{A,\omega}
 \qquad(\omega\in\Omega,\ x\in A_\omega),
\]
then $\widehat\mu(\overline A)\leq\widehat\mu(\overline B)$.
\end{lemma}
\begin{proof}
The nonzero-map slope inequality gives
\[
 \widehat\mu_{\min}(\overline A)
 \leq\widehat\mu_{\max}(\overline B)+h(f),
\]
where $h(f)$ is the extended map height obtained from the logarithms of the
completed operator norms. The local contractions give $\|f_\omega\|\leq1$,
so monotonicity of the lower integral gives $h(f)\leq0$.
Semistability and Hermitianity of the source give
$\widehat\mu_{\min}(\overline A)=\widehat\mu(\overline A)$;
semistability of the target gives
$\widehat\mu_{\max}(\overline B)=\widehat\mu(\overline B)$.
Substitution proves the assertion. The two mean slopes are finite real
numbers, so their comparison in the extended real line is the same comparison
in $\mathbb R$.
\end{proof}

This is the nonzero-map form of the comparison in
\cite[Proposition 2.4 and the proof of Theorem 6.2]{bost2026}.
Only the source's minimal-slope characterization requires Hermitianity.
In particular, the comparison does not require the map to be injective.

For integers $a,k\geq1$, let $L_{a,ak}$ denote the integral rectangular
Specht lattice with its word Gram metric, and put
\[
 C_a(k)=\frac{\log\operatorname{covol}(L_{a,ak})}
                  {\operatorname{rank}_{\mathbb Z}L_{a,ak}}.
\]
Thus $C_a(k)=c_a(ak)$ in the notation of this paper.

\begin{lemma}[Eliminating rectangular Specht errors]
\label{lem:specht-error-limit}
\source{bost-adelic-2026:page-0028}
\uses{lem:specht-limit}
\lean{BostConjecture.Adelic.le_of_rectangleLogCovolume_bounds}
\leanok
Let $I$ be a finite set, let $s,d\geq1$ and $a_i\geq1$ for $i\in I$,
and let $l_r,k_{r,i}$ be nonnegative integers for $r\geq0$ such that
\[
 n_r:=d(r+1)=s l_r=a_i k_{r,i}\qquad(i\in I).
\]
If $c,x,y\in\mathbb R$ and, for every $r\geq0$,
\[
 n_r x-cC_s(l_r)\leq n_r y+c\sum_{i\in I}C_{a_i}(k_{r,i}),
\]
then $x\leq y$. There is no sign restriction on $c$.
\end{lemma}
\begin{proof}
The common degree $n_r$ is positive, so every width appearing in the formula
is positive. Since $n_r\to\infty$ and the heights are fixed and positive,
$l_r=n_r/s\to\infty$ and $k_{r,i}=n_r/a_i\to\infty$.
The normalized covolume limit therefore gives
\[
 \frac{C_s(l_r)}{s l_r}\longrightarrow0,
 \qquad
 \frac{C_{a_i}(k_{r,i})}{a_i k_{r,i}}\longrightarrow0
 \quad(i\in I).
\]
Divide the assumed inequality by $n_r>0$, and use the equalities of degrees:
\[
 x-c\frac{C_s(l_r)}{s l_r}
 \leq y+c\sum_{i\in I}\frac{C_{a_i}(k_{r,i})}{a_i k_{r,i}}.
\]
A finite sum preserves convergence, as does multiplication by the fixed real
number $c$. The two sides tend to $x$ and $y$, respectively. The order relation
on $\mathbb R$ is closed, hence $x\leq y$.
\end{proof}

\section{The auxiliary bundles and the finite-degree comparison}

From now on $K$ is countable of characteristic zero.
For a positive-rank Hermitian bundle $\overline A$ of rank $a$ and a positive
multiple $n=ak$, write
\[
 S_{a,n}(\overline A)
   =(\det\overline A)^{\otimes k}\otimes\overline L^S_{a,n}.
\]
Here $\overline L^S_{a,n}$ is the constant adelic bundle of the integral Specht
lattice: it has the word Gram norm at infinity and the integral-basis maximum
norm at finite places. The determinant line carries the actual determinant
norm. These choices fix both the module and every completed norm in the
auxiliary bundle.

\begin{lemma}[Slope of the determinant power]
\label{lem:determinant-power-slope}
\uses{def:adelic-slope,def:adelic-induced-norms}
\source{bost-adelic-2026:page-0025,bost-adelic-2026:page-0026}
\lean{BostConjecture.Adelic.determinantPower_slope}
\leanok
For an adelic bundle $\overline A$ and $k\geq0$, the orthogonal tensor power
of its determinant line has slope
\[
 \widehat\mu((\det\overline A)^{\otimes k})
      =k\widehat{\deg}(\overline A).
\]
The empty power is the scalar line of degree zero.
\end{lemma}
\begin{proof}
The determinant bundle is a line, and every tensor power of lines has rank
one, including the empty power. Its slope is therefore its line degree.
Additivity of line degree under tensor products gives $k$ times the line
degree of $\det\overline A$. By definition this last line degree is
$\widehat{\deg}(\overline A)$. For $k=0$, the sum of line degrees is empty
and equals zero.
\end{proof}

Let $I$ be a finite nonempty index set and let $(\overline E_i)_{i\in I}$ be
positive-rank Hermitian bundles. Give
$W=\bigotimes_{i\in I}E_i$ its orthogonal tensor norms, and let
$\iota:U\hookrightarrow W$ be a nonzero subspace with induced norms.
Set $s=\dim_K U$ and $a_i=\dim_K E_i$. For a positive common multiple $n$ of
$s$ and all $a_i$, put
\[
 B_n=S_{s,n}(\overline U),\qquad
 D_n=\bigotimes_{i\in I}S_{a_i,n}(\overline E_i^{\vee}).
\]
The canonical determinant maps give
$\iota_U:B_n\to U^{\otimes n}$ and
$j_n:D_n\to(W^{\vee})^{\otimes n}$. Define the original-field pairing
\[
 \Psi_n:B_n\longrightarrow D_n^{\vee},\qquad
 \Psi_n(b)(d)=\langle j_n(d),\iota^{\otimes n}(\iota_U(b))\rangle.
\]
The brackets are algebraic tensor evaluation; no averaging operator is part
of this definition.

\begin{lemma}[The finite-degree geometric slope bound]
\label{lem:finite-degree-geometric-slope}
\source{bost-adelic-2026:page-0026,bost-adelic-2026:page-0027}
\uses{lem:semistable-contraction-slope,lem:determinant-power-slope,lem:specht-degree,lem:specht-semistability,cor:group-line-twist-semistability,lem:independent-group-semistability,lem:slope-duality,lem:local-pairing-contraction}
\notready
In the preceding setting, if $\Psi_n\neq0$, then
\[
 n\widehat\mu(\overline U)-m_\infty C_s(n/s)
 \leq n\sum_{i\in I}\widehat\mu(\overline E_i)
       +m_\infty\sum_{i\in I}C_{a_i}(n/a_i).
\]
\end{lemma}
\begin{proof}
The permutation action on each rectangular Specht lattice is absolutely
irreducible and preserves every completed lattice norm. Give each determinant
line power the trivial action. Tensoring the lattice representation with this
one-dimensional trivial representation preserves absolute irreducibility:
after any field extension, choosing a nonzero vector of the line identifies
the resulting representation with the lattice representation. The action
preserves the tensor norm because it preserves both factor norms. All these
bundles are Hermitian.

Consequently $B_n$ is semistable. Each factor of $D_n$ has its own
permutation group action, so the independent-group semistability theorem
makes $D_n$ and $D_n^{\vee}$ semistable. This assertion uses the product of
the groups, not their diagonal subgroup.
The local pairing contraction gives $\|\Psi_{n,\omega}\|\leq1$ at every
place. Since $\Psi_n\neq0$, the semistable contraction comparison yields
\[
 \widehat\mu(B_n)\leq\widehat\mu(D_n^{\vee}).
\]

The lattice degree formula and the determinant-power formula give
\[
 \widehat\mu(B_n)=n\widehat\mu(\overline U)-m_\infty C_s(n/s).
\]
For each $i$, the same formulas and Hermitian duality give
\[
 \widehat\mu(S_{a_i,n}(\overline E_i^{\vee}))
   =-n\widehat\mu(\overline E_i)-m_\infty C_{a_i}(n/a_i).
\]
Repeated use of the binary tensor-slope formula computes the slope of $D_n$
as the sum of these expressions. Duality negates that sum. Substituting both
computed slopes in the comparison proves the displayed inequality.
\end{proof}

\section{Geometric semistability and quotient norms}

The Pl\"ucker point of $U\subset W$ is the line
$\bigwedge^s U\subset\bigwedge^s W$.
It is geometrically semistable if, over an algebraic closure of $K$, a
positive-degree homogeneous polynomial invariant under
$\prod_{i\in I}\operatorname{SL}(E_i)$ is nonzero on a nonzero vector of
that line.

\begin{theorem}[Geometric subspace slope inequality]
\label{thm:geometric-subspace-bound}
\source{bost-adelic-2026:page-0026,bost-adelic-2026:page-0027,bost-adelic-2026:page-0028}
\uses{lem:nonzero-pairings,lem:finite-degree-geometric-slope,lem:specht-error-limit}
\notready
If the Pl\"ucker point of the nonzero subspace $\iota:U\hookrightarrow W$
is geometrically semistable, then
\[
 \widehat\mu(\overline U)\leq\sum_{i\in I}\widehat\mu(\overline E_i).
\]
No arithmetic semistability hypothesis is imposed on the $\overline E_i$.
\end{theorem}
\begin{proof}
The nonzero-pairing theorem supplies $d\geq1$ divisible by $s$ and all $a_i$
such that $\Psi_{d(r+1)}\neq0$ for every integer $r\geq0$.
Put $l_r=d(r+1)/s$ and $k_{r,i}=d(r+1)/a_i$; divisibility makes these positive
integers. The finite-degree bound applies at every one of these degrees.
Apply the error-elimination lemma with
\[
 c=m_\infty,\qquad x=\widehat\mu(\overline U),\qquad
 y=\sum_{i\in I}\widehat\mu(\overline E_i).
\]
Its conclusion is the required inequality.
\end{proof}

\begin{lemma}[Dual norm of an orthogonal tensor quotient]
\label{lem:tensor-quotient-dual-norm}
\uses{def:adelic-induced-norms}
\source{bost-adelic-2026:page-0028}
\lean{ArakelovGeometry.NormFamily.orthogonalTensor_pushforward_dual_norm,ArakelovGeometry.NormFamily.orthogonalTensorPair_pushforward_dual_norm}
\leanok
Let $q:W\twoheadrightarrow Q$ be a surjective linear map from the orthogonal
tensor product of a finite family of Hermitian normed spaces over $S$.
Give $Q$ the pushforward norms. Under the canonical tensor-duality isomorphism,
let
\[
 j:Q^{\vee}\xrightarrow{q^{\vee}}W^{\vee}
       \xrightarrow{\ \sim\ }\bigotimes_{i\in I}E_i^{\vee}.
\]
At every place and on every completed vector $\lambda\in Q_\omega^{\vee}$,
\[
 \|\lambda\|_{Q^{\vee},\omega}
   =\|j_\omega(\lambda)\|_{\bigotimes_i E_i^{\vee},\omega}.
\]
The same statement holds for the ordinary binary tensor product.
\end{lemma}
\begin{proof}
Fix a place. Write $C=\|q^{\vee}(\lambda)\|_{W^{\vee}}$.
For every lift $x\in W_\omega$ of $y\in Q_\omega$,
$|\lambda(y)|\leq C\|x\|$. Taking the infimum over all lifts gives
$|\lambda(y)|\leq C\|y\|_Q$, hence $\|\lambda\|_{Q^{\vee}}\leq C$.
Conversely $\|q(x)\|_Q\leq\|x\|$, so
$|\lambda(q(x))|\leq\|\lambda\|_{Q^{\vee}}\|x\|$ for every $x$.
Thus $C\leq\|\lambda\|_{Q^{\vee}}$. This proves equality for the transpose.
The canonical tensor-duality isomorphism is isometric for the prescribed
Hermitian tensor norms, so composing with it gives the asserted identity.
Scalar extension commutes with these canonical maps, hence the computation
applies to all completed vectors. The binary statement uses the binary
tensor-duality isomorphism and the identical argument.
\end{proof}

\begin{lemma}[Binary tensor norms under finite indexing]
\label{lem:binary-tensor-norm}
\uses{def:adelic-induced-norms}
\lean{ArakelovGeometry.NormFamily.orthogonalTensor_norm_boolTensorEquiv}
\leanok
Let $\overline E,\overline F$ be Hermitian normed spaces over $S$. Set
$A_0=E$ and $A_1=F$. The canonical map
\[
 \beta:\bigotimes_{i\in\{0,1\}}A_i\longrightarrow E\otimes_K F,
 \qquad \bigotimes_i x_i\longmapsto x_0\otimes x_1,
\]
preserves the orthogonal tensor norm at every place, on every completed vector.
\end{lemma}
\begin{proof}
Identify the two-element index set with the disjoint union of two singletons,
in the indicated order. Reindexing preserves the Hilbert tensor norm because
it permutes the orthonormal tensor basis, and preserves the epsilon tensor
norm because it bijects the families of norm-bounded dual test vectors.
The binary tensor norm is the finite tensor norm transported through the
canonical grouping of those two singletons. On pure tensors the composite
of reindexing and grouping sends $\bigotimes_i x_i$ to $x_0\otimes x_1$;
multilinearity identifies it with $\beta$ on the whole tensor product.
These maps commute with scalar extension, so the same argument applies in
each completion and proves the claimed norm identity.
\end{proof}

\begin{theorem}[Geometric quotient slope inequality]
\label{thm:geometric-quotient-bound}
\source{bost-adelic-2026:page-0028}
\uses{thm:geometric-subspace-bound,lem:geometric-dual-quotient,lem:tensor-quotient-dual-norm,lem:binary-tensor-norm,lem:slope-duality}
\notready
Let $\overline E$ and $\overline F$ be positive-rank Hermitian adelic bundles
over the countable characteristic-zero field $K$, and put
$\overline T=\overline E\otimes\overline F$ with its orthogonal tensor norms.
Let $q:E\otimes_K F\twoheadrightarrow Q$ be a positive-rank quotient and
let $\overline B=q_*\overline T$ be its actual pushforward bundle.
If $q$ is geometrically semistable, meaning that its dual Pl\"ucker point
is geometrically semistable for the contragredient tensor action, then
\[
 \widehat\mu(\overline E)+\widehat\mu(\overline F)
     \leq\widehat\mu(\overline B).
\]
Neither factor is required to be arithmetically semistable.
\end{theorem}
\begin{proof}
The dual-quotient bridge identifies the geometric hypothesis with geometric
semistability of
\[
 j:Q^{\vee}\hookrightarrow E^{\vee}\otimes_K F^{\vee}
\]
for the full special-linear groups of the dual factors. Surjectivity of $q$
implies injectivity of its transpose, and tensor duality is an isomorphism,
so $j$ is injective. Its source is nonzero because $\dim_K Q>0$.
The binary norm identity identifies this target with the two-factor instance
of the finite-family tensor bundle used in the subspace theorem.
The dual-quotient norm identity says that the induced norms along $j$ are
precisely those of $\overline B^{\vee}$. Apply the geometric subspace
inequality to the Hermitian bundles $\overline E^{\vee}$ and
$\overline F^{\vee}$. It yields
\[
 \widehat\mu(\overline B^{\vee})
 \leq\widehat\mu(\overline E^{\vee})+
       \widehat\mu(\overline F^{\vee}).
\]
The tensor bundle and its quotient $\overline B$ are Hermitian. Hermitian
slope duality therefore negates all three slopes in this inequality.
Multiplication by $-1$ gives the assertion, with the original pushforward
norms on $\overline B$.
\end{proof}
